% !TEX root = ../thesis-main.tex

\RequirePackage{mathtools}  % mathtools asks to be loaded before unicode-math...
\documentclass{kthpq-thesis}

% ==== PACKAGES ====

% ---- kthpq-thesis ----

% Custom biblatex package to have cool macros
\usepackage[%
  bibstyle=numeric,
  maxnames=99,
  isbn=false,
  url=false,
  urldate=ymd,
  defernumbers=true
]{biblatex-kthpq}

% \makeatletter  % bibliography numbers in margin...
% \defbibenvironment{bibliography}
% {\list{\printtext[labelnumberwidth]{%
%     \printfield{prefixnumber}%
%     \printfield{labelnumber}}}
% {\setlength{\bibhang}{0pt}%
% \setlength{\leftmargin}{\bibhang}%
% \setlength{\itemindent}{-\leftmargin}%
% \setlength{\itemsep}{\bibitemsep}%
% \setlength{\parsep}{\bibparsep}}}
% {\endlist}
% {\item}
% \makeatother

% ---- Formatting ----

\usepackage{parskip}

\usepackage[
  % demo  % omit inclusion of PDFs
]{pdfpages}

% longtable: ignoring this error https://github.com/latex3/latex2e/issues/1907
% with VS Code LaTeX Workshop 
\usepackage{longtable}

% ---- Math ----

% mathtools loaded at the top!
\usepackage{amsthm,thmtools}  % thmtools for QED in any theorem

% ---- Floats ----

% \usepackage{flafter}
\usepackage{caption}
\usepackage{subcaption}

% ---- Other ----

\usepackage{tikz}
\usetikzlibrary{babel,calc,cd}
\usepackage[ruled]{algorithm2e}

\usepackage{imakeidx}

% ---- Load last ----

% Hyperref needs to come after parskip
% (https://tex.stackexchange.com/questions/450551/theorems-and-parskip#comment1853746_450551),
% and generally should be loaded last.
\usepackage[
%  draft, % Uncomment to remove all links (useful for printing in black and white)
%  hidelinks,
  % bookmarksdepth,% w/o argument: uses tocdepth, supposedly
  pdfusetitle,
  colorlinks,
  linktoc=all,
  allcolors={digital blue}
]{hyperref}

% amsmath (from unicode-math) > hyperref > cleveref
\usepackage[noabbrev,capitalize]{cleveref}

% ==== COMMANDS ====

% \WarningFilter{latex}{Marginpar on page}
\DeclarePrintbibliographyDefaults{title=References}

% ---- Font ----

\setmainfont{Georgia Pro}[
  Uprightfont = *-Regular,
  BoldFont = *-Bold,
  ItalicFont = *-Italic,
  BoldItalicFont = *-BoldItalic,
  Numbers = OldStyle,
  Scale = \fontscale
]

% ---- Proofing ----

% \geometry{showframe}  % layout frames
% \usepackage{lua-visual-debug}  % more boxes!
% \overfullrule=10pt  % visual indicator of overfull boxes

% ---- Bibliography ----

\addbibresource{bib/mine.bib}
\addbibresource{bib/kappa.bib}
\addbibresource{bib/CI.bib}

\DeclareFieldFormat{eprint:github}{%
GitHub: \href{https://www.github.com/\strfield{eprint}}{\texttt{\strfield{eprint}}}
}

% ---- Layout ----

% Custom header and footer with fancyhdr
\fancypagestyle{fancy}{%
  \fancyhead{}  % reset
  \fancyhead[LO]{\sffamily\small\rightmark}
  \fancyhead[RE]{\sffamily\small\leftmark}
  \fancyhead[LE,RO]{\sffamily\small\thepage}  % Left for even pages, right for odd pages
}

\fancypagestyle{plain}{% for special pages, e.g. first pages of chapters
  \fancyhead{}  % reset
  \fancyhead[LE,RO]{\sffamily\small\thepage}  % Left for even pages, right for odd pages
}

\makeatletter
% Flipbook stuff: inspired by nilsvu/latex-flipbook on GitHub
% Current hunch is that pgfpicture sets the bounding box based on the vertices (the last path?)
\newif\if@frontmatter
\fancyfoot{}  % reset footer
\fancyfoot[RO]{%
  \if@mainmatter%
    \ifnum\thepage<60\relax%
      \AddToHookNext{shipout/background}{\put(\paperwidth-30mm,3mm-\paperheight){\input{figs/flip/leo-\fpeval{trunc((\thepage - 1) / 2, 0)}.pgf}}}%
    \fi%
  \else%
    \if@frontmatter%
      \AddToHookNext{shipout/background}{\put(\paperwidth-30mm,3mm-\paperheight){\input{figs/flip/leo-0.pgf}}}%
    \fi%
  \fi%
}

% Remove fancy headers from front matter and add check for front matter
\let\oldfrontmatter\frontmatter
\let\oldmainmatter\mainmatter
\let\oldbackmatter\backmatter
\def\frontmatter{\oldfrontmatter\@frontmattertrue\pagestyle{plain}}
\def\mainmatter{\oldmainmatter\@frontmatterfalse\pagestyle{fancy}}
\def\backmatter{\oldbackmatter\@frontmatterfalse}
\makeatother

\setcounter{tocdepth}{1}

% Try to make things more compact
%\linespread{.9}

% ---- Sections & titles ----

\newcommand{\inleftmargin}[2][\marginparsep]{%
  \makebox[0pt][r]{#2\hspace{#1}}%
}

\titleformat{\section}
  {\sffamily\fontsize{13pt}{13pt}\selectfont\bfseries}
  {\inleftmargin{\thesection}}{0em}{}

\titleformat{\subsection}
  {\sffamily\fontsize{12pt}{12pt}\selectfont\bfseries}
  {\inleftmargin{\thesubsection}}{0em}{}

% New section-like for papers
\makeatletter
\newcommand\papersection[2][]{% OPTIONAL: shorttitle, misctext
  \setkeys{kthpq@papersection}{#1}%
  \@papersection{#2}%
  \let\@misctext\defaultmisctext%
  \let\@shorttitle\@nil
}

\let\@shorttitle\@nil
\define@key{kthpq@papersection}{misctext}{\def\@misctext{#1}}
\define@key{kthpq@papersection}{shorttitle}{\def\@shorttitle{#1}}

\DeclareCiteCommand{\@papersection}
{%
  \pagetrackerfalse%
  \citetrackerfalse%
  \usebibmacro{prenote}%
  \defcounter{maxnames}{99}%
  \sffamily%
  \let\theoldsection\thesection
  \renewcommand{\thesection}{\papername\ \Alph{section}}
}
{%
  \ifx\@shorttitle\@nil
    \section{\thefield{labeltitle}}
  \else
    \section[\@shorttitle]{\thefield{labeltitle}}
  \fi
  % % automatically add label (outside labels don't detect this section)
  % \label{PS:\strfield{entrykey}}
  \printnames{author}.
  \ifentrytype{article}{%
    In \usebibmacro{tryfjournal},
  }{}%
  \ifentrytype{inbook}{%
    In \printfield{booktitle},
  }{}%
  \ifentrytype{incollection}{%
    In \printfield{booktitle},
  }{}%
  \ifentrytype{inproceedings}{%
    In \printfield{booktitle},
  }{}%
  \ifentrytype{misc}{%
    \@misctext\ifx\@misctext\empty\else, \fi%
  }{}%
  \printfield{year}.
  \printfield{doi}\printfield[eprint:arxiv]{eprint}%
}
{}
{%
  \usebibmacro{postnote}%
  \let\thesection\theoldsection
}

% Adding bibliographic data to ToC
% can't patch this command for some reason so here's the full thing
\renewcommand{\papertitlepage}[2][]{%
  \cleardoublepage%
  \newgeometry{top=30mm,bottom=30mm,left=43mm,right=\@tabmargin}%
  \thispagestyle{empty}%
  \setkeys{kthpq@papertitlepage}{#1}
  \phantomsection
  \paper@tab% before stepping counter for better positioning
  \refstepcounter{paper}
  \@papertitlepage{#2}
  \renewcommand{\thechapter}{\thepaper}
  \markboth{}{}% Reset header
  \addtocontents{toc}{%  <-- inserted block here
    \protect\def\protect\@misctext{\@misctext}%
    \protect\@paperdetailsintoc{#2}%
  }% end inserted block
  \let\@misctext\defaultmisctext
  \restoregeometry
  \vfil\newpage\null
  \thispagestyle{empty}
  \afterpage{\nopagecolor}% Reset the color after the current page
  \newpage
}

\def\@paperdetailsintoc#1{%
  \hspace{\cftchapnumwidth}%
  \parbox[t]{\dimexpr\textwidth - \cftchapnumwidth - \@tocrmarg\relax}{%
    \sffamily\@paperdetailsintoc@i{#1}
  }%
  \vspace{5pt}% manually adjusted
}

\DeclareCiteCommand{\@paperdetailsintoc@i}
  {%
    \pagetrackerfalse%
    \citetrackerfalse%
    \defcounter{maxnames}{99}%
  }
  {%
    \printnames{author}.
    \ifentrytype{article}{%
      In \usebibmacro{tryfjournal},
    }{}%
    \ifentrytype{inbook}{%
      In \printfield{booktitle},
    }{}%
    \ifentrytype{incollection}{%
      In \printfield{booktitle},
    }{}%
    \ifentrytype{inproceedings}{%
      In \printfield{booktitle},
    }{}%
    \ifentrytype{misc}{%
      \@misctext\ifx\@misctext\empty\else, \fi%
    }{}%
    \printfield{year}.
    \printfield{doi}\printfield[eprint:arxiv]{eprint}%
  }
  {}
  {}
\makeatother

% ---- Theorems ----

\declaretheoremstyle[
  headfont=\sffamily\bfseries,
  headformat={\inleftmargin[.5\marginparsep]{\mdseries\NUMBER}\NAME \NOTE},
  headpunct={.\,},
  bodyfont={\itshape}
]{theorem-margin}

\declaretheorem[style=theorem-margin, parent=chapter]{theorem}
\declaretheorem[style=theorem-margin, numbered=no, name=Theorem]{theorem*}
\declaretheorem[style=theorem-margin, sibling=theorem]{proposition}
\declaretheorem[style=theorem-margin, sibling=theorem]{corollary}
\declaretheorem[style=theorem-margin, sibling=theorem]{lemma}

\declaretheoremstyle[
  headfont=\sffamily\bfseries,
  headformat={\inleftmargin[.5\marginparsep]{\mdseries\NUMBER}\NAME \NOTE},
  headpunct={.\,},
  qed=$\mdlgblksquare$
]{plain-margin}

\declaretheorem[style=plain-margin, sibling=theorem]{definition}
\declaretheorem[style=plain-margin, sibling=theorem]{remark}
\declaretheorem[style=plain-margin, sibling=theorem]{example}
\declaretheorem[style=plain-margin, sibling=theorem]{notation}

% \declaretheorem[style=plain-bird, numbered=no, name=Definition]{definition**}

% Add a footnote right before the period of the theorem head.
% This version is designed to be used with cleveref. The commented out lines (1)
% should work if cleveref is not loaded. With cleveref, there are two types of
% theorems, those that are siblings and those that are not (I think?). In any
% case, it calls one of the two following command: \cref@thmoptarg and
% \cref@thmnoarg (run \meaning\@thm to see how this is done). We locally
% redefine both of them just to be sure.
\makeatletter
\newenvironment{footnoted}[2]{%
  \def\@currentthmname{#1}
  % \let\old@thm\@thm% (1)
  % \def\@thm##1##2##3{\old@thm{##1\thm@headpunct{\footnotemark.}}{##2}{##3}}% (1)
  \let\old@cref@thmoptarg\cref@thmoptarg
  \def\cref@thmoptarg[##1]##2##3##4{\old@cref@thmoptarg[##1]{##2\thm@headpunct{\footnotemark.}}{##3}{##4}}
  \let\old@cref@thmnoarg\cref@thmnoarg
  \def\cref@thmnoarg##1##2##3{\old@cref@thmnoarg{##1\thm@headpunct{\footnotemark.}}{##2}{##3}}
  \csname#1\endcsname
  \footnotetext{#2}
}{%
  \csname end\@currentthmname\endcsname
}
\makeatother

% ---- floats ----

\captionsetup[subfigure]{justification=centering}

% ---- non-math commands ----

% epigraph
\makeatletter
\newenvironment{epigraph}[1][\@nil]{%
  \def\epigraph@uthor{#1}
  \begin{flushright}
    \small \itshape
    \begin{tabular}{l@{}}
}{%
    \end{tabular}%
    \vspace{1ex}
  \ifx\epigraph@uthor\@nnil
  \else
    \\
    \normalfont
    --- \epigraph@uthor
  \fi
  \end{flushright}%
  \vspace{1em}%
}
\makeatother

\newcommand{\ir}[1]{{\color{brick} #1}}
\def\mydinkus{\par\noindent\hfill\textasciitilde~$\cdot$~\textasciitilde\hfill\null\par}

% in-margin highlights
\makeatletter
\def\notate{\@dblarg\kthpq@highlight}
\def\term{\@dblarg\kthpq@term}
\newcommand{\kthpq@term}[2][]{\kthpq@highlight[#1]{\textbf{#2}}}

% % Fix value (3) for same-note line spacing,
% % Fix value (2) for vertical alignment with the text.
% % After fixing (2) and (3), fix value (1) for multiple-note spacing.
% \marginparpush=1.2\baselineskip  % (1)
% \newcommand{\kthpq@highlight}[2][]{%
%   \leavevmode% I don't remember why
%   \checkoddpage
%   \marginpar{%
%     \ifoddpageoroneside
%       \raggedright
%     \else
%       \raggedleft
%     \fi
%     \vspace{-5pt}  % (2)
%     \begin{spacing}{.7}  % (3)
%       \scriptsize\sffamily #1
%     \end{spacing}
%   }%
%   \penalty\@M% no linebreak after hspace
%   \hspace{0pt}% allow hyphenation by inserting a zero-width thing
%   #2%
% }
% \makeatother

% \newcommand{\fallbackhighlight}[1]{%
%   \marginnote{%
%     \ifoddpageoroneside%
%       \raggedright%
%     \else%
%       \raggedleft%
%     \fi%
%     \scriptsize\sffamily #1%
%   }
% }

% % Uncomment to remove margin notes
\makeatletter
\newcommand{\kthpq@highlight}[2][]{\index{#1}#2}
\makeatother
\newcommand{\fallbackhighlight}[1]{\index{#1}}
\makeindex[intoc]